\pfl{}

This has not been a pleasant book to research or to write. It will not be
a pleasant book to read. It deals with a profoundly disturbing topic: a
system of justice, designed to protect children from sexual abuse, that
has gone so terribly wrong that it often abets the very evils it exists to
correct. The failure of any judicial system is bad enough; the failure of
the American family court 1 system, resulting in the systematic mishan-
dling of sex abuse charges, exposes children throughout the country to
grave dangers. That is something anyone would rather not know, given
the choice. \pfl{}


But the truth is simply too serious to ignore. If it is true—as we have
found—that the American family court system, built to safeguard the
welfare of children and their families, has all too often become a place
where civil rights are thrust aside, where the search for truth is subor-
dinated to “junk science,” and where children are forced into unwanted
relationships with parents they have accused of sexually abusing them—
while the parents who try to protect them are punished with the loss of
custody—then none of us should be indifferent to family court malfunc-
tion. While the family courts treat too many mothers who believe that
their children have been sexually abused as guilty until proven inno-
cent (guilty of parental unfitness, of hysteria, of “hostility” toward the
accused parent), and take away their children as a result, ignorance of
the system is clearly too high a price for any of us to pay. \pfl{}


Indeed, the American public is already aware that something is wrong
with the system that is supposed to protect our children. Over the last
\notePage{lxxiii}
twenty years, the child welfare system has attracted increasing pub-
lic attention and controversy. Again and again, the public has learned
of astonishing judicial lapses as a result of which children are kept in
the custody of physically or sexually abusive parents—until it is too
late. \pfl{}


Yet there has always been much more to the story than this. The public
is aware that the family court system sometimes fails to protect children
from sexual abuse. But it is only just becoming aware that the system
may actually place children with sexually abusive parents because the
other parent (usually the mother) is stigmatized for making the accu-
sation in the first place. That is the phenomenon we examine in this
book. \pfl{}


Although ours is the first book to analyze this pattern in objective de-
tail, we are not the first to notice the evils we examine. Well over a decade
ago, several researchers, including activist attorney H. Joan Pennington,
coined the term “protective parent” to describe those parents, usually
mothers, who engage the family court system to protect their children
from alleged sexual abuse by the child’s father. In many cases, Pen-
nington complained, these parents faced a fierce backlash from judges,
law guardians, caseworkers, court-appointed mental health experts, and
others. 2 \pfl{}


Pennington (now retired) founded the National Center for Protec-
tive Parents, based in Trenton, New Jersey, in 1990. Over the following
six years, the Center published a number of well-documented studies
on the plight of mothers in the family courts. In 1992, Pennington tes-
tified before the House Judiciary Committee on how to educate judges
about the protective parent problem. 3 Pennington has argued that fam-
ily courts not only fail, in many cases, to stop the abuse of children but
actually go much further by punishing mothers for trying to protect
them. \pfl{}


In 1999, Kim Gandy, then vice-president of the National Organi-
zation for Women (NOW), was moved to write to one of us that this
problem cried out for legislative reform: “[T]here is a frightening trend
in both the juvenile and family law courts of the United States, which is
placing literally thousands of children in danger. In court cases across
\notePage{lxxiv}
the county, children who have disclosed abuse by a parent . . . are being
placed in the sole custody of the abusive parent Congress must be made
aware of how prevalent this problem is.” \pfl{}


In April 2004, Kathryn Mazierski, president of the NOW chapter
in New York State, echoed Gandy’s denunciation of the family courts.
Mazierski, joined by other NOW state chapters and national child ad-
vocacy organizations, demanded a federal investigation of the family
courts. Mazierski’s press release did not mince words: “in case after
case, the mother who makes the report of sex abuse is punished—
stripped by the court of her custodial rights. This phenomenon can
be observed in so many states, and involves such profound viola-
tions of the rights of mothers and the children concerned, that a
federal investigation into the Constitutional and civil rights violations
of mothers caught in the family court system is the best way to seek a
remedy.” \pfl{}


According to Mazierski, “One out of every two calls I get on a daily
basis is from a mother who just lost custody because she tried to get the
family court to protect her child from sexual abuse committed by the
father.” Similarly, Rita Henley Jensen, editor in chief of Women’s eNews,
an online daily news service reaching twenty thousand members of the
press, has told one of us that “the bulk of the daily incoming emails to
Women’s eNews are desperate pleas from mothers who have lost custody
of their children they tried in vain to protect from sexual abuse.” 4 This
suggests that, if anything, the plight of protective parents is worsening:
six years earlier, the director of the Violence Against Women Office of
the US Justice Department had written to one of us that at least “two calls
a week” received by her office were from protective parents “desperate
for assistance.” 5 \pfl{}


If this sounds to you like the product of judicial madness, you are
not alone. During a 1997 Mother’s Day press conference held on the
steps of the Capitol in Washington, DC, a devastated Georgia mother,
whose contact with her five-year-old son had been completely sus-
pended by a family court, shouted: “Family court is a bad dream that
you can’t wake up from!” The other fifty mothers at the rally plainly
agreed. \pfl{}



\notePage{lxxv}
Our own look into the “bad dream” of the family court system be-
gan in 1986, when Dr. Amy Neustein, one of the authors of this book,
lost custody of her six-year-old daughter to a man an eyewitness had
accused of sexually molesting her—despite the confirming account of
the girl herself and the corroboration of experts. Appalled by what she
had experienced, Dr. Neustein founded the Help Us Regain The Chil-
dren (HURT) research center in Brooklyn, New York, in 1986, moving
to New York City in 1999. \pfl{}


HURT advertised itself to protective mothers primarily by word of
mouth, as well as through local and national television talk show appear-
ances, during which the Center’s address and phone number would be
posted on the screen. Referrals were also made to HURT by women’s or-
ganizations in New York and nationwide, and battered women’s shelters
listed the Center’s phone number as a resource in their printed literature. \pfl{}


The results were astonishing. HURT was contacted in well over four
thousand cases. About one thousand met the study requirements set
forth by HURT, which required its research subjects to provide as full
a court record as possible: motion papers, court transcripts, decisions,
and orders. Later, this book’s other author, Michael Lesher, Esq., un-
dertook a legal review of HURT’s cases and added others from his
own experience. The authors found that those cases tragically con-
firmed the warnings of Pennington, Gandy, and Mazierski. Family
courts really were ignoring evidence of sexual abuse and awarding cus-
tody to fathers against whom credible cases of sexual abuse had been
presented. \pfl{}


Of those mothers who lost custody to allegedly abusive ex-husbands
or former partners, the majority had been primary caretakers of their
children—the main nurturers—and in a number of cases the children
had been very young when they were removed from their mothers. Some
had not even been weaned. Others had disabilities that made them par-
ticularly dependent on their mothers. None of these facts made any
consistent difference in the results of the cases. Healthy children and
sick ones, teenagers and toddlers, all were just as likely to be removed
from their mothers when allegations emerged of sexual abuse by the
father. \pfl{}



\notePage{lxxvi}
We found that such results cut straight across the American demo-
graphic pie. Both rich and poor women lost custody of their children
to alleged sex offenders, as did women at different educational levels.
The judicial outcomes were largely the same in large metropolises and
in rural areas. \pfl{}


The evidence we have collected has enabled us to demonstrate in
detail—and to analyze—the backlash against protective parents that
activists have complained about for decades. We offer this book as a
detailed analysis of a system that has turned far too much of its energy
toward the punishment, not of parents who abuse their children, but
of those parents (almost always women) who reasonably suspect their
spouse or partner of child sexual abuse. \pfl{}


Just as the evidence we have amassed enables us to remove any doubt
as to the reality of this problem, it has enabled us to address two impor-
tant reservations sometimes voiced by skeptics. These should be dealt
with at once. \pfl{}


First, we do not claim—nor do we believe—that the miscarriages of
justice we dissect in this book characterize all family court cases across
the country. But it remains true, unfortunately, that severely mishan-
dled cases are not uncommon. As we will show, other investigations,
formal and informal, have found—as we did—that far too many cases of
alleged sexual abuse have ended with inexplicable punishments meted
out to a mother whose “crime” was to try to protect her child from
sexual abuse. And the consistency of this finding is even more strik-
ing because court records of sex abuse cases are nearly always barred
from public review, making large-scale surveys unusually difficult. The
same may be said of the sheer volume of desperate calls and e-mails
from bewildered protective mothers to women’s organizations and other
groups. \pfl{}


Take the case of Childhelp USA. Founded in 1959 as a groundbreak-
ing purveyor of child abuse prevention information, Childhelp later
established a telephone hotline for child abuse “crisis counseling,” the
number of which (1-800-4-a-child) slowly began to make appearances
on daytime television in the 1980s. By 1992, after several years of run-
ning the hotline, Childhelp officials had received so many calls from
\notePage{lxxvii}
mothers claiming that they had tried to protect their children from
abuse, only to have the children taken from their custody, that they
decided to investigate the phenomenon more fully. Dubbed the “Sys-
tem Failure Project,” Childhelp’s investigation was straightforward: a
detailed questionnaire was mailed to a several thousand callers who
had complained of losing custody of a child after reporting sexual abuse
by the other parent. The questionnaire addressed the specific actions
of the salient figures in each case: judges, law guardians, social ser-
vice caseworkers, court-appointed mental health experts, and others.
A full 10 percent of the questionnaires were completed and returned
(a large figure for any survey, particularly one involving several pages
of questions). The results showed such consistent accusations of mis-
conduct, from hundreds of independent respondents, that Childhelp
decided it was necessary to address legislators about the seriousness
of the problem. Mariam Bell, Childhelp’s director of governmental af-
fairs at the time, approached one congressional committee after another,
telling lawmakers what her organization had found and urging them to
convene a federal hearing into the issue. 6 \pfl{}


Thus, whatever the exact proportion of family court cases involv-
ing a “backlash” against protective mothers, it is impossible to doubt
that the problem is serious, widespread, and intractable. Indeed, in this
book, we have attempted to go beyond numbers to identify the socio-
logical and legal patterns that result in such outcomes—thus revealing
an entrenched system that will continue to wrong allegedly abused chil-
dren, and their mothers, until it is checked. It is difficult to say exactly
how large, numerically, this problem is. But no one can question its
gravity. \pfl{}


Second, we are well aware that not every father accused of sexual
abuse is actually guilty. What is more, we know that as the courts
operate—and must operate—even parents who are guilty of sexual
abuse must sometimes remain uncorrected because of the absence of
acceptable proof. We do not seek a system in which the accused is de-
nied the presumption of innocence. Nevertheless, the all-too-common
voices that complain of an assault against innocent men in the courts,
that worry loudly that sex abuse charges pose a lopsided threat to
\notePage{lxxviii}
the accused—with the inevitable corollary that such charges must be
received with skepticism—are seriously wrong. Public fears of a “witch
hunt” against fathers are not borne out by the facts. The existing lit-
erature clearly establishes that false or malicious accusations of sexual
abuse, even in custody litigation, are not at all common. Dziech and
Schudson (1989) point to several sources—Thoennes, Pearson, and
Tjaden (1988), Quinn (1988), Berliner (1988), and Corwin, Berliner,
Goodman, Goodwin, and White (1987)—as examples of well-respected
studies showing that “even in family courts, false allegations [of sexual
abuse] remain rare. . . . [M]ost . . . are substantiated” (p. 204). More
recent empirical studies, as well as critical analyses of the published re-
search, confirm these findings (Denike, Huang, and Kachuk 1998; Faller,
Corwin, and Olafson 1993; Faller and DeVoe 1995; McDonald 1998;
Thoennes and Tjaden 1990). \pfl{}


Nor is it true that sex abuse charges must be met with judicial resis-
tance because such accusations are all too easily made and all too readily
endorsed by credulous psychologists. Again, the professional literature
belies this claim. The exact opposite seems more likely to be the case:
experts often fail to diagnose abuse when it has actually occurred. In a
1992 article in the Journal of Interpersonal Violence, psychologists Law-
son and Chaffin stressed the difficulty of obtaining a disclosure from an
abused child. Lawson and Chaffin interviewed twenty-eight children,
ranging from three-year-olds to adolescents, and found that more than
half (57 percent) of those with sexually transmitted diseases did not re-
port having been sexually abused during the initial interview. 7 What is
more, even when credible evidence supports a sex abuse charge, there
is no reason to believe it will help the legal position of the mother who
makes it. For example, an ominous study by Faller and DeVoe (1995) at
the University of Michigan found that forty mothers concerned about
possible sexual abuse of their children had incurred “negative sanctions”
from a court of law, which included being jailed, losing custody (to the
alleged offender, to a relative, or to foster parents), losing some or all
visitation rights, being ordered not to report alleged abuse again to the
court (or to child protective services or the police), and being barred
from taking the child to a doctor or a therapist because of sex abuse
\notePage{lxxix}
concerns. Disturbingly, the researchers found that the parents who had
received such punitive measures had actually scored higher on a com-
posite scale of likelihood of sexual abuse—a calculation of the likelihood
that their child had been abused, taking into account such factors as the
existence of medical evidence—than those parents who had not received
such “sanctions.” 8 \pfl{}


It follows that the sufferings of protective parents in family courts
cannot be explained away with pat answers. To understand what is hap-
pening, one must delve into the culture of the family courts, carefully
documenting what occurs, and then explaining how such acts fit within
a framework that makes sense to the system’s practitioners—although
to anyone else they look like sheer madness. \pfl{}


That is the purpose of this book. \pfl{}


As the reader will find, there is no shortage of madness to discuss.
Family courts can punish protective mothers in a bewildering variety of
ways. They can, and often do, transfer custody of the allegedly abused
child to the parent accused of abusing him or her. In fact, some judges go
further, forbidding mothers to see their children more than a few hours
a month, and even then, only under strict supervision. In some cases, a
mother’s visits have been suspended indefinitely because she suspected
sexual abuse. Judges have been known to ban telephone calls, letters,
and e-mails between mother and child. In one California case, a mother
was criticized by a family court judge for waving to her children when
she saw them in another car on the highway. \pfl{}


Hand-in-hand with family court madness is a growing pattern of
parental mutiny—a refusal by protective parents to accept the proce-
dures or orders of courts whose logic they can no longer comprehend.
Faced with judicial rulings that seem to them contrary to the most ba-
sic parental instincts and values, they have begun to rebel. In defiance
of family court orders that they believe put their children at risk, they
litigate furiously, take their stories to the press, go on hunger strikes,
or turn fugitive with their children with the help of a modern-day
“Underground Railroad.” \pfl{}


Like judicial madness, parental rebellion against the family courts has
taken many forms. Protective mothers have organized grassroots groups
\notePage{lxxx}
throughout the country—in New York, Ohio, Michigan, Georgia, North
Carolina, Washington, and California. The groups monitor court cases,
hold candlelight vigils, and in some cases actively support the Under-
ground Railroad that hides mothers who have gone on the run with their
children. \pfl{}


Some mothers have pursued legislative reform, becoming unpaid ac-
tivists. They offer lawmakers and administrators a grim look into the
family courts. In response, a few legislators have become champions of
protective mother causes. Some have actually drafted legislation on the
subject, although little of it has reached the statute books. In his book, In
Defense of Public Order (1961), political theorist Harold Lasswell has ar-
gued that “the history of legislation often shows that pressure groups . . .
are developed long before an issue becomes a burning question on the
calendar of the Congress or state legislatures.” 9 This is clearly true for
the fighting protective parents. \pfl{}


Our goal in this book is to provide a thorough picture of the infes-
tation of the family courts by a special form of judicial madness, and
to illustrate its causes and its consequences. We will show how family
court judges and law guardians, child welfare workers, court-appointed
mental health experts, visitation supervisors, and others have created
a system that all too often betrays its own principles. We will also dis-
cuss some practical reforms that can help to redirect the system toward
healthier goals. \pfl{}


It should be stressed that our study of the backlash against moth-
ers by the family courts is an empirical analysis, concentrating on how
more than why. That is, we have avoided speculation as to the underlying
psychological or historical forces that may cause family court personnel
to react so violently when mothers accuse fathers of child sexual abuse.
This is not to say that some writers, aware of the same phenomenon, have
not attempted this. One attempt in particular, the theory put forth by the
feminist writer Louise Armstrong, deserves mention here as probably
the best known. \pfl{}


Armstrong argues that charges of sexual abuse provoke such a hostile
response from family courts because they pose a threat to the inher-
ent patriarchal structure of the family. In Rocking the Cradle of Sexual
\notePage{lxxxi}
Politics: What Happened When Women Said Incest (1994), Armstrong
writes that fathers accused of sexual abuse have answered such charges
with “a political war, a war to safeguard their rights.” And she adds that
men are winning the war: charges of sexual abuse have left women “vul-
nerable to charges of hysteria or vindictiveness,” whereas “in men sheer
rage can be taken as reasonable outrage, for righteous indignation: for a
sign that their circumstance is indeed demanding of remedy” (p. 138). \pfl{}


Such a theory is, of course, extremely difficult to prove. It is clear that
women’s greater aggressiveness in domestic litigation has inspired the
creation of “fathers’ rights” groups (the activities of which are noted
in Chapter 1) and that some of these organizations have agitated for
a return to “biblical” patriarchal rights. But this extreme is a fringe
phenomenon. Our research has shown that male judges and social work-
ers have no monopoly on hostility to sex abuse charges that women
level at men. The same backlash can be observed whether the family
court personnel are predominantly male or female, liberal or conser-
vative, religious or agnostic. We have chosen, therefore, to concentrate
on methods that dissect family court madness as an institutional phe-
nomenon, describing what the system clearly is—we are less concerned
with inevitably more debatable attempts to reveal the madness’s ultimate
source. \pfl{}


A word about our research methods is certainly in order here. As
previously noted, it is impossible to undertake a large-scale, random
review of family court case files because these are not available to the
public. However, we have been as thorough as circumstances permit.
The research for this book spans twenty years and states all across
the United States. Cases were analyzed longitudinally over all or most
of their litigation history, including appellate review when applicable.
Other accessible and relevant sources were also examined: reports of
statewide judicial conduct commissions and gender-bias task forces,
published case digests of decisions, and testimonial records of state and
federal public hearings. We interviewed judges, private attorneys, law
guardians, social workers, and mental health experts. Litigants were in-
terviewed on multiple occasions over extended periods of time. On a
few occasions, we were able to observe videotapes of trials. \pfl{}



\notePage{lxxxii}
To enhance this book’s analysis of family courts, we have employed
a form of sociological inquiry known as ethnomethodology. Using this
approach, we sought to discover, describe, and analyze the collabora-
tive ways in which “social interactants” in the family courts accomplish
the situated production of social order in their day-to-day activities. For
example, this book attempts to explain how judges, as members of a fam-
ily court “culture,” methodically 10 re-create, on each occasion that they
rule on the issues before them, the family court entity itself as a socially
organized system. In developing this analysis, we reviewed the existing
ethnomethodological studies of courts and court processes (Atkinson
1979; Atkinson and Drew 1979; Drew 1992; Holstein 1983, 1988; Hol-
stein and Gubrium 1995; Komter 1995, 1998; Lynch 1997; Matoesian
1993; Maynard 1984; Meehan 1997; Pollner 1979; Travers and Manzo
1997; Watson 1990, 1997), although none of these was directed at the
study of the family court institution per se. \pfl{}


One particularly important feature of our case research is our em-
phasis on studying the progress of a single case over time. Judicial
madness often follows an evolutionary process. As a bitterly contested
case progresses, rulings may become increasingly bizarre. As that hap-
pens, madness may sometimes be glimpsed merely by reading the court
record as it unfolds over time. For example, a Georgia family court judge,
whose ruling on a child neglect question that had caused two young chil-
dren to be placed in foster care was reversed on appeal as “not supported
by the appropriate considerations and findings of fact,” 11 reportedly in-
sisted to the litigants that the appellate judges’ order would not affect
him, saying, “There’s nothing they can do to me; I’m bullet proof!” \pfl{}


Parental “mutiny,” likewise, is an evolutionary process. Some parents
who have never before participated in a political protest begin to chal-
lenge the system after being caught up in it and then develop, step by
step, into ardent activists. Some parents choose a more radical path.
After protracted court proceedings, in which (they believe) issues of
child safety are being dangerously mishandled, these parents defy court
orders outright. They become either fugitives or prisoners. \pfl{}


One other comment seems appropriate here. As we have written, this
book is an objective study, presenting objective data. But objectivity
\notePage{lxxxiii}
is not the same as indifference. Throughout this book, we have been
at pains to be accurate, but when the facts we reveal are outrageous,
we have not minced words about them. Nor do we expect our read-
ers to be dispassionate about what we describe, for the truth about the
family courts gives the lie to some cherished American assumptions.
Americans believe they are safe from arbitrary abuse of governmental
power—yet child protective services, “law guardians,” and family court
judges can cast aside the norms of due process when suppressing an alle-
gation of sexual abuse. Americans believe that constitutional protections
follow them everywhere—but family courts are an exception, at least for
protective parents; as one family court judge put it, “There is no First
Amendment in my court.” Americans believe that sexist stereotypes
have been soundly expelled from American political life—yet family
courts still stigmatize mothers who make abuse allegations as “hysteri-
cal,” “vindictive,” and “hostile,” as if, in family courts, the appearance of
an angry woman still scandalizes. We see no point in trying to minimize
the injustice of such phenomena. \pfl{}


On the other hand, simply to identify the abuses perpetrated by the
family court system is not enough. One must also try to understand how
they arise, how they become part of an institutional phenomenon—and
we have devoted considerable effort to this question. Finally, one must
consider how the system can be reformed. And we address that issue as
well. \pfl{}


This book is divided into three parts. The first gives an overview of
family court malfunction and the parental mutiny that results from it.
We outline the madness that has infected the family courts and describe
the new legal landscape that has helped make the madness possible. We
show how the system has reacted—or rather, failed to react—to severe
criticism from commentators, newspapers, and even legislators. We also
explain the roles that judges, lawyers, and other official “auxiliaries” play
in the family court process. \pfl{}


For example, one important modern development is the increased
role of Child Protective Services (CPS) agencies, whose workers have
formulated their own criteria for determining who is a fit parent. In
many cases, CPS workers have charged mothers with child neglect
\notePage{lxxxiv}
solely on the strength of a claim that these mothers “brainwashed” their
children into believing they had been sexually abused. \pfl{}


Corresponding to this new legal landscape is a new sort of parental
rebellion, and we describe this, too, including the formation of a new
“Underground Railroad” for mothers of allegedly abused children and
the women who have risked their freedom to lead it. \pfl{}


The book’s second part presents our observations in greater depth.
First, we discuss in some detail the history and applications of eth-
nomethodologists’ approach to sociological inquiry and demonstrate
how this approach can be applied to the study of the family courts.
Next, we present detailed observations of the abuse of judicial powers
in family courts—both the kinds of abuse and reasons they seem to oc-
cur. We discuss the function of “law guardians,” lawyers intended to
advocate solely for the welfare of children in abuse cases who all too
often have other agendas. We discuss in depth the behavior of child
protective services agencies and how they, too, often fail to protect the
children they exist to protect. We also chart the evolution of contro-
versial mental health theories that have been given great evidentiary
weight in the family courts. These theories—with names such as Ma-
licious Mother Syndrome (MMS) and Parental Alienation Syndrome
(PAS)—are unique to family litigation and seem to have been created
specifically for use in counterattacking abuse allegations in court. Fi-
nally, we treat in detail the behavior of mothers who rebel in various
ways against the system and the effects on them of the courts’ violent
disruption of their family life. \pfl{}


The third, and last, part of this book discusses possible ways of
revamping the system. In so doing, we examine what is required to “re-
birth” the system, starting with a radical change in the beliefs that drive
it. In addition, we examine a number of legislative proposals that have
been introduced in New York, California, and Texas with the aim of im-
proving the family courts. We add to these existing proposals a number
of suggested reforms to redress particular shortcomings of the system
that we describe in this book. Our suggestions are aimed not only at
judges but also at law guardians and social workers, who are an integral
part of the system. \pfl{}



\notePage{lxxxv}
The full significance and impact of the failures and abuses this book
describes will only be clearly seen many years from now. But even now it
is possible to document how a court system designed to protect families
has been tearing too many children and innocent mothers apart.
